Persistent memory lets a mobile personal agent maintain preferences, plans, and safety-relevant facts across sessions, but it also converts transient disclosures into a durable attack surface. Existing Android search primitives authorize at package or schema granularity, while contemporary agent-memory systems optimize what to retain and retrieve. Neither boundary expresses that a memory acquired for one purpose may be consumed only by particular apps and tools, until a specified expiry, nor does it conservatively govern summaries derived from multiple sources. We present \ms, a purpose-bound memory service for mobile personal agents. MemScope places a reference monitor between Android agents and a system-owned persistent store. Every item binds content to provenance, authorized purposes and consumers, app/tool identities, sensitivity, and lifetime; every read carries an authenticated retrieval token derived from the active AppFunction/task context. Authorization precedes semantic ranking, so rejected items do not enter the model context. Derived memories inherit the meet of their cited sources' labels and are admitted only after source and attribution checks. Retrieval-time expiry and background compaction enforce lifetime, while policy-aware eviction separates safety-critical retention from ordinary relevance. Under explicit trust, complete-mediation, and token-authenticity assumptions, we prove admission soundness, monotone restriction of derived labels, returned-plaintext cross-scope non-disclosure, and expiry safety; these properties deliberately do not claim semantic non-inference after an authorized disclosure. We specify an AOSP implementation and a preregistered evaluation on Pixel and Samsung phones using published memory-extraction and poisoning attacks, long-term-memory benchmarks, and multi-app workloads. Because the project contains no completed measurements, all curves and result tables are marked as conceptual or pending rather than presented as empirical evidence.
